% !TEX program = xelatex
% 大物实验 · 实验 3.8 · 原理与流程精讲（lecture-deepdive 规范）
\documentclass[11pt]{ctexart}
%=====================================================================
% theme-violet.tex —— 讲义精讲 · 紫花主题
% 用法：\documentclass[11pt]{ctexart}
%       \input{assets/theme-violet}
%       \begin{document} ... \end{document}
% 只用 ctexart 自带能力 + xcolor/enumitem/fancyhdr/pifont/amsmath，
% 不依赖 tikz/tcolorbox，保证 MiKTeX 可直接编译。
%=====================================================================

\usepackage[a4paper,left=2.0cm,right=2.0cm,top=2.1cm,bottom=2.3cm,%
            headsep=12pt,footskip=22pt]{geometry}
\usepackage{amsmath,amssymb}
\usepackage{xcolor}
\usepackage{array}
\usepackage{enumitem}
\usepackage{pifont}
\usepackage{fancyhdr}
\usepackage[hidelinks]{hyperref}

%---------------------------------------------------------------- 配色
\definecolor{violet}{RGB}{107,78,143}      % 主色：紫
\definecolor{violetdeep}{RGB}{63,44,92}    % 深紫：标题/强调
\definecolor{petal}{RGB}{183,159,208}      % 花瓣：装饰线、花饰
\definecolor{bandviolet}{RGB}{247,243,251} % 浅紫底：档案/提示框
\definecolor{bandcream}{RGB}{252,249,242}  % 米白底：易错/注意框
\definecolor{rulepurple}{RGB}{217,207,230} % 分隔线
\definecolor{inkgray}{RGB}{92,96,110}      % 次要文字

%---------------------------------------------------------------- 字体
% 有 Noto CJK 用 Noto，没有就退回 ctex 默认（Windows 下为宋体/黑体）
\IfFontExistsTF{Noto Serif CJK SC}{\setCJKmainfont{Noto Serif CJK SC}}{}
\IfFontExistsTF{Noto Sans CJK SC}{\setCJKsansfont{Noto Sans CJK SC}}{}

%---------------------------------------------------------------- 版式
\setlist{nosep,leftmargin=1.9em}
\linespread{1.28}
\setlength{\parskip}{3.5pt}
\setlength{\parindent}{0pt}
\setlength{\tabcolsep}{5pt}
\renewcommand{\arraystretch}{1.28}

%---------------------------------------------------------------- 花饰
% \ding{95} = 五瓣小花，\ding{96} = 星形花；都来自 pifont，必定存在
\newcommand{\orn}{\textcolor{petal}{\ding{95}}}          % 行内小标
\newcommand{\ornstar}{\textcolor{petal}{\ding{96}}}      % 星形花
\newcommand{\flowerrule}{%                                % 居中花饰分隔线
  \par\vspace{7pt}%
  \noindent\makebox[\linewidth]{\color{petal}%
    \rule{0.20\linewidth}{0.6pt}\hspace{8pt}{\small\ding{95}}\hspace{4pt}%
    \raisebox{1.1pt}{\tiny\ding{95}}\hspace{4pt}{\small\ding{95}}\hspace{8pt}%
    \rule{0.20\linewidth}{0.6pt}}%
  \par\vspace{7pt}}

%---------------------------------------------------------------- 标题与节头
\newcommand{\doctitle}[3]{%   #1 主标题  #2 副标题  #3 元信息（可用 \\ 分行）
  \begin{center}
    {\color{petal}\small \orn\enspace #2\enspace \ornstar}\\[7pt]
    {\LARGE\bfseries\color{violetdeep} #1}\\[10pt]
    {\color{petal}\rule{0.55\linewidth}{1pt}}\\[4pt]
    {\small\color{inkgray} #3}
  \end{center}\vspace{2pt}}

\newcommand{\hdr}[2]{%   #1 章节号（如 §0）  #2 标题
  \par\vspace{12pt}%
  \noindent\colorbox{violetdeep}{\color{white}\bfseries\small\;#1\;}\hspace{7pt}%
  {\color{violetdeep}\large\bfseries #2}\par\nobreak
  \vspace{3pt}{\color{petal}\hrule height 1.2pt}\vspace{6pt}}

\newcommand{\subhdr}[1]{%
  \par\vspace{6pt}%
  \noindent{\color{violet}\bfseries #1}\par\vspace{1pt}}

%---------------------------------------------------------------- 提示框
\newcommand{\callout}[1]{% 浅紫底：讲义档案、读法说明
  \par\medskip\noindent
  \fcolorbox{rulepurple}{bandviolet}{\parbox{\dimexpr\linewidth-2\fboxsep-2\fboxrule}{\vspace{2pt}#1\vspace{2pt}}}%
  \par\medskip}

\newcommand{\warnbox}[1]{% 米白底：易错、边界条件
  \par\medskip\noindent
  \fcolorbox{rulepurple}{bandcream}{\parbox{\dimexpr\linewidth-2\fboxsep-2\fboxrule}{\vspace{2pt}#1\vspace{2pt}}}%
  \par\medskip}

\newcommand{\gapitem}[2]{% #1 缺口类型（紫） #2 内容
  \par\vspace{2pt}\noindent
  \hangindent=1.4em\hangafter=1
  {\color{petal}\small\ding{96}}\enspace{\bfseries\color{violetdeep}#1}\enspace #2\par}

% 前提旁注：只放在「需要理解的那一步」后面，给陈述 + 本题代入，不重证。
\newcommand{\aside}[1]{%
  \par\vspace{1pt}\noindent
  \hangindent=1.6em\hangafter=1
  {\small\color{inkgray}%
    \hspace*{0.6em}{\color{petal}$\triangleright$}\hspace{2pt}#1}%
  \par\vspace{2pt}}

%---------------------------------------------------------------- 表格
% 本机 MiKTeX 的 colortbl 与 array 版本不匹配（\rowcolor 会报错），
% 故表格统一用「上粗紫线 + 表头 + 中细线」的三线式，不依赖 colortbl。
\newcommand{\tabtop}{\par\vspace{3pt}{\color{violet}\hrule height 1.2pt}\vspace{1pt}}
\newcommand{\tabmid}{\par\vspace{1pt}{\color{petal}\hrule height 0.5pt}\vspace{1pt}}
\newcommand{\tabbot}{\par\vspace{3pt}{\color{rulepurple}\hrule height 0.5pt}\vspace{8pt}}
\newcommand{\thead}[1]{{\bfseries\color{violetdeep}#1}}

%---------------------------------------------------------------- 页眉页脚
\pagestyle{fancy}
\fancyhf{}
\renewcommand{\headrulewidth}{0.4pt}
\renewcommand{\headrule}{\hbox to\headwidth{\color{rulepurple}\leaders\hrule height \headrulewidth\hfill}}
\fancyhead[L]{\small\color{inkgray}\songti \doctag}
\fancyhead[R]{\small\color{inkgray}\songti \coursetag}
\fancyfoot[C]{\small\color{inkgray}\orn\enspace{\color{violet}\thepage}\enspace\orn}

% 每篇讲义在导言区设这两个值
\newcommand{\doctag}{}
\newcommand{\coursetag}{}
\newcommand{\settags}[2]{\renewcommand{\doctag}{#1}\renewcommand{\coursetag}{#2}}

%---------------------------------------------------------------- 数学记号
\newcommand{\vect}[1]{\boldsymbol{#1}}
\newcommand{\dd}{\,\mathrm{d}}
\newcommand{\ddx}{\frac{\mathrm{d}}{\mathrm{d}x}}
\addtolength{\jot}{2pt}



\settags{实验 3.8 原理与流程精讲}{大物实验}
\usepackage{tikz}
\usepackage{graphicx}
\usetikzlibrary{arrows.meta,calc,positioning}
\usepackage{booktabs}
\usepackage{float}

\tikzset{
  comp/.style={fill=white,draw=black,thick},
  arr/.style={->,>=Stealth,thick},
  lab/.style={font=\small},
  dimlab/.style={font=\footnotesize\color{inkgray}}
}

\begin{document}

\doctitle{铁磁材料的磁化曲线与磁滞回线}
{实验 3.8 原理与流程精讲 $\cdot$ 从磁畴到示波器上的 B-H 回线}
{依据：Experiment 3.8 讲义（9 页）与 Preview Requirements（3 页）$\cdot$ 2026-10-31（周六）20:00--22:00 $\cdot$ 综合楼 207 $\cdot$ 2026-10-27 整理}

%=====================================================================
\hdr{§0}{本讲坐标、记号与报告任务}

本讲义把两份材料合成一条线：先讲清\textbf{铁磁性为什么不是线性的}（§1--§2），再讲\textbf{磁滞回线怎么读}（§3），然后讲\textbf{示波器为什么能把 \(B\) 和 \(H\) 画成 \(x\) 和 \(y\)}（§4），最后落到\textbf{进实验室怎么做、数据怎么算}（§5--§6）。§7 与 §8 分别是 10 道预习问题和 8 道课后思考题的简答，§9 是易错点，§10 是两份材料里没有明说、本讲义补上的缺口清单。

\callout{\textbf{两份材料的页码约定}：本讲义中的「讲义 pX」指 \texttt{Experiment3.8.MagnetizationCurve\-andHysteresisLoop.pdf} 的物理页（共 9 页）；「预习 pX」指 \texttt{Experiment3.8previewquestions.pdf} 的物理页（共 3 页）。实验日期按预约截图：2026-10-31（周六）20:00--22:00，综合楼 207，任课教师王泽，3 学时。}

\subhdr{§0.1 记号表（全文统一）}

\begin{center}
\begin{tabular}{@{}p{1.7cm}p{9.0cm}p{3.2cm}@{}}
\toprule
\thead{符号} & \thead{含义} & \thead{单位} \\
\midrule
\(B\) & 磁感应强度（磁通密度） & T = Wb/m\(^2\) \\
\(H\) & 磁场强度 & A/m \\
\(M\) & 磁化强度 & A/m \\
\(\mu_0\) & 真空磁导率，\(4\pi\times10^{-7}\) & T\(\cdot\)m/A = H/m \\
\(\mu_r\) & 相对磁导率，\(\mu_r=\mu/\mu_0\) & 无量纲 \\
\(N_1,N_2\) & 原边、副边匝数 & 匝 \\
\(L\) & 磁路平均长度（不是样品「长度」随便取） & m \\
\(S\) & 铁芯截面积 & m\(^2\) \\
\(R_1,R_2\) & 原边取样电阻、副边积分电阻 & \(\Omega\) \\
\(C\) & 积分电容 & F \\
\(S_x,S_y\) & 示波器 CH1、CH2 的垂直灵敏度 & V/div \\
\(x,y\) & 屏幕上相对中心的格数 & div \\
\(B_r\) & 剩磁（residual magnetism） & T \\
\(H_c\) & 矫顽力（coercive force） & A/m \\
\(B_s\) & 饱和磁感应强度 & T \\
\bottomrule
\end{tabular}
\end{center}

单位换算只需记一条：由 \(1\ \mathrm{F}=1\ \mathrm{A\cdot s/V}\)、\(1\ \Omega=1\ \mathrm{V/A}\)，得
\[
1\ \mathrm{V}\cdot\mathrm{F}\cdot\Omega/\mathrm{m^2}
=1\ \mathrm{V\cdot s/m^2}
=1\ \mathrm{Wb/m^2}
=1\ \mathrm{T}.
\]
预习 p2 的附录写的 \(1\ \mathrm{T}=1\ \mathrm{V}\cdot\mathrm{F}\cdot\Omega/\mathrm{m^2}\) 就是这条；后面把 \(C R_2S_y y/(N_2S)\) 算成 T 时，用的也是它。

\subhdr{§0.2 报告要交什么}

\callout{\textbf{预习报告}：10 道预习问题（中英文答案，精简）；\textbf{实验报告}：实验目的、实验器材、课后思考题、约 50 词实验体会；\textbf{数据表}：自己打印带去，记录 \(R_1,R_2,C,S_x,S_y\)，样品 1 与样品 2 的 \(B_r,H_c\)，样品 1 的基本磁化曲线 \(x\)--\(y\)--\(y/x\) 与最大相对磁导率。}

%=====================================================================
\hdr{§1}{铁磁性从哪里来：磁畴}

\subhdr{1.1 原子磁矩与磁畴}

铁磁材料（铁、钴、镍及其合金）的磁性来自原子尺度的磁矩。相邻原子的磁矩在交换相互作用下倾向于平行排列，于是在小区域（约微米量级）内自发地取同一方向，这样的区域叫\textbf{磁畴}（domain）。一块未磁化的铁磁体内部有许多磁畴，相邻磁畴的磁矩方向不同，宏观上互相抵消，所以整块材料不显磁性（图 1；原讲义 Figure 3.8-1 见下方原图）。

加外磁场后，磁畴的变化分三段：

\begin{enumerate}
\item \textbf{畴壁移动（可逆）}：与外场方向接近的磁畴长大，方向不利的磁畴缩小。小磁场下撤去外场，畴壁能弹回，这一段是可逆的。
\item \textbf{畴壁移动（不可逆）}：外场继续增大，畴壁被材料中的缺陷、杂质、晶界「钉」住，只能跳跃式前进（巴克豪森跳变）。撤去外场后畴壁回不到原位置，这就是\textbf{磁滞}的微观来源，也是能量损耗的来源。
\item \textbf{磁畴转动}：所有磁畴逐渐转到外场方向；再增大外场，磁化强度 \(M\) 趋于饱和值 \(M_s\)，磁感应强度 \(B\) 只能靠 \(\mu_0H\) 这一项继续缓慢增加。
\end{enumerate}

\begin{figure}[htbp]
\centering
\begin{tikzpicture}[scale=0.78]
% left: unmagnetized
\draw[thick] (0,0) rectangle (4.4,4.4);
\draw[arr] (0.7,1.0)--(1.3,0.7);
\draw[arr] (2.6,0.8)--(2.2,1.4);
\draw[arr] (3.4,1.1)--(3.8,1.7);
\draw[arr] (1.0,2.4)--(0.6,3.0);
\draw[arr] (2.2,2.6)--(2.9,2.9);
\draw[arr] (3.6,2.5)--(3.1,3.2);
\draw[arr] (1.4,3.6)--(2.0,3.9);
\node[dimlab] at (2.2,-0.45) {未磁化：各畴磁矩互相抵消};
% right: magnetized
\begin{scope}[xshift=6.2cm]
\draw[thick] (0,0) rectangle (4.4,4.4);
\foreach \y in {0.7,1.5,2.3,3.1,3.9} {
  \draw[arr] (0.5,\y)--(1.5,\y);
  \draw[arr] (2.0,\y)--(3.0,\y);
  \draw[arr] (3.5,\y)--(4.0,\y);
}
\node[dimlab] at (2.2,-0.45) {外场中：磁畴转向外场方向};
\end{scope}
\end{tikzpicture}
\caption{磁畴示意（重绘，依据原讲义 p2 的文字描述；箭头只表示磁矩方向，不表示大小；原图见下）。}
\end{figure}

\begin{center}
\includegraphics[width=0.55\linewidth]{fig-3.8-orig/fig3.8-1.png}\\[2pt]
{\small\color{inkgray} 原讲义 Figure 3.8-1（裁剪）：左为未磁化的随机磁畴，右为磁化后磁矩取向一致。}
\end{center}

\subhdr{1.2 居里温度与软磁/硬磁}

温度升高到\textbf{居里温度} \(T_C\) 以上，热运动破坏交换相互作用，磁畴瓦解，材料变成顺磁体。讲义 p1 给的例子是磁铁矿（magnetite）\(T_C=580^\circ\mathrm{C}\)；常见金属的 \(T_C\) 约为铁 \(770^\circ\mathrm{C}\)、钴 \(1115^\circ\mathrm{C}\)、镍 \(358^\circ\mathrm{C}\)（这三个数是教材常识，不是讲义原文）。

按矫顽力大小，铁磁材料分两类：

\begin{itemize}
\item \textbf{软磁}（soft magnetic）：矫顽力小、磁导率高、回线窄，容易磁化也容易退磁，用于变压器、继电器、电机铁芯（硅钢、软铁）。
\item \textbf{硬磁}（hard magnetic）：矫顽力大、剩磁高、回线宽，撤去外场后仍保留强磁性，用于永磁体（钕铁硼、钐钴、锶铁氧体）。
\end{itemize}

%=====================================================================
\hdr{§2}{\(B\)、\(H\)、\(M\) 与磁化曲线}

\subhdr{2.1 三个量的关系}

磁感应强度 \(B\)、磁场强度 \(H\)、磁化强度 \(M\) 在 SI 制中满足
\[
B=\mu_0(H+M).
\]
真空中 \(M=0\)，所以 \(B=\mu_0H\)。对线性磁介质，\(M=\chi H\)，于是
\[
B=\mu_0(1+\chi)H=\mu_0\mu_rH,\qquad \mu_r=1+\chi.
\]
对铁磁材料，\(M\) 与 \(H\) 的关系不仅非线性，还与磁化历史有关；这时「磁导率」只能逐点定义：
\[
\mu=\frac{B}{H},\qquad \mu_r=\frac{B}{\mu_0H}.
\]
注意这里的 \(\mu=B/H\) 是\textbf{割线磁导率}（secant permeability），不是 \(B\)--\(H\) 曲线的切线斜率 \(\dd B/\dd H\)（微分磁导率）。讲义 p3 说「the slope of the magnetization curve at any given point gives the permeability」，这句话把两者混用了；本实验报告里按预习 p3 的公式用 \(\mu_r=B/(\mu_0H)\)，也就是割线值。详见 §10 的 G3。

\subhdr{2.2 磁化曲线：从退磁态开始}

把样品完全退磁（\(B=H=0\)，记作 O 点），然后\textbf{单调增大} \(H\)，测出的 \(B\)--\(H\) 曲线叫\textbf{基本磁化曲线}（magnetization curve）。它的形状分四段：

\begin{enumerate}
\item \textbf{初始段}：\(H\) 很小，畴壁可逆移动，\(B\) 增加较慢，\(\mu\) 较小；
\item \textbf{陡峭段}：畴壁不可逆跳跃，\(B\) 增加最快，\(\mu\) 达到最大值 \(\mu_{\max}\)；
\item \textbf{膝部}：畴壁移动接近完成，磁畴开始转动，\(B\) 的增长变慢；
\item \textbf{饱和段}：\(M\to M_s\)，再增大 \(H\) 只有 \(\mu_0H\) 的贡献，\(B\approx\mu_0(H+M_s)\)，曲线几乎成为斜率 \(\mu_0\) 的直线，对应的 \(B_s\) 叫饱和磁感应强度。
\end{enumerate}

\begin{figure}[htbp]
\centering
\begin{minipage}{0.49\linewidth}
\centering
\begin{tikzpicture}[x=0.72cm,y=0.72cm]
\draw[->] (0,0)--(5.7,0) node[right,font=\small]{\(H\)};
\draw[->] (0,0)--(0,4.7) node[above,font=\small]{\(B\)};
\draw[thick,inkgray] plot[domain=0:5.3,smooth,variable=\x] (\x,{0.42*\x});
\draw[thick,violet] plot[domain=0:5.3,smooth,variable=\x] (\x,{4*\x/(\x+0.8)});
\draw[thick,petal] plot[domain=0:5.3,smooth,variable=\x] (\x,{3*\x/(\x+1.5)});
\node[dimlab,right] at (4.4,1.55) {空气};
\node[dimlab,right] at (2.55,3.75) {铁};
\node[dimlab,right] at (4.1,2.75) {钢};
\end{tikzpicture}
\par\small (a) 空气、铁、钢的 \(B\)--\(H\) 曲线（示意）
\end{minipage}\hfill
\begin{minipage}{0.49\linewidth}
\centering
\begin{tikzpicture}[x=0.72cm,y=0.72cm]
\draw[->] (0,0)--(5.7,0) node[right,font=\small]{\(H\)};
\draw[->] (0,0)--(0,4.7) node[above,font=\small]{\(\mu_r/10^3\)};
\draw[thick,violet] plot[domain=0.05:5.3,smooth,variable=\x] (\x,{2+4*\x/(1+\x*\x)});
\draw[dashed,inkgray] (0.9,0)--(0.9,4.1);
\node[dimlab,right] at (1.0,4.2) {\(\mu_{\max}\)};
\end{tikzpicture}
\par\small (b) 相对磁导率随 \(H\) 的变化（示意）
\end{minipage}
\caption{磁化曲线与磁导率曲线（重绘，依据原讲义 Figure 3.8-2、3.8-3；曲线为示意，不是实测数据；原图见下）。}
\end{figure}

\begin{center}
\begin{minipage}{0.49\linewidth}\centering
\includegraphics[width=\linewidth]{fig-3.8-orig/fig3.8-2.png}\\[1pt]
{\footnotesize\color{inkgray} 原 Figure 3.8-2：空气、铁、钢的 B--H 曲线}
\end{minipage}\hfill
\begin{minipage}{0.49\linewidth}\centering
\includegraphics[width=\linewidth]{fig-3.8-orig/fig3.8-3.png}\\[1pt]
{\footnotesize\color{inkgray} 原 Figure 3.8-3：$\mu_r$ 随 $H$ 的变化}
\end{minipage}
\end{center}

\subhdr{2.3 为什么 \(B\)--\(H\) 不是直线}

把 \(B=\mu_0(H+M)\) 拆成两项看：\(\mu_0H\) 是「真空也会有的那一份」，随 \(H\) 线性增长；\(\mu_0M\) 是材料贡献的那一份，由磁畴状态决定。磁畴的畴壁移动和转动不是 \(H\) 的线性函数：小磁场时只有少数畴壁在动，\(M\) 增长慢；中等磁场时大量畴壁跳跃，\(M\) 增长快；高磁场时所有磁矩已经排齐，\(M\) 饱和到 \(M_s\)，再增大 \(H\) 只剩 \(\mu_0H\)。所以整条 \(B\)--\(H\) 曲线先缓、后陡、再平，\(\mu=B/H\) 自然不是常数。

%=====================================================================
\hdr{§3}{磁滞回线：\(B_r\)、\(H_c\) 与损耗}

\subhdr{3.1 一个完整循环}

从退磁态 O 出发，把激励电流（也就是 \(H\)）从 0 单调增大到正最大值，\(B\) 沿初始磁化曲线上升到 a 点（接近饱和）；然后把电流单调减到 0，\(B\) 不沿原路返回，而是降到 b 点：此时 \(H=0\)，但 \(B=B_r\neq0\)，这个 \(B_r\) 叫\textbf{剩磁}。要把它降到 0，必须反向通电到 c 点，对应的反向磁场强度大小叫\textbf{矫顽力} \(H_c\)。继续反向增大电流到 d 点（反向饱和），再减到 0 到 e 点（反向剩磁 \(-B_r\)），再正向增大到 f 点（反向矫顽力），最后回到 a 点。a--b--c--d--e--f--a 这条闭合曲线就是\textbf{磁滞回线}（图 3；原讲义 Figure 3.8-4 见下方原图）。

\begin{figure}[htbp]
\centering
\begin{tikzpicture}[x=1.0cm,y=1.0cm]
\draw[->] (-4.8,0)--(4.8,0) node[right]{\(H\)};
\draw[->] (0,-4.8)--(0,4.8) node[above]{\(B\)};
\draw[thick,violet,domain=-4:4,smooth,samples=120] plot (\x,{4*tanh((\x+0.8)/1.5)});
\draw[thick,violet,domain=-4:4,smooth,samples=120] plot (\x,{4*tanh((\x-0.8)/1.5)});
\draw[dashed,inkgray,domain=0:4,smooth] plot (\x,{4*\x/(\x+0.8)});
\fill (4,4) circle (2.2pt) node[above right,font=\small]{\(a\)};
\fill (0,{4*tanh(0.8/1.5)}) circle (2.2pt) node[above right,font=\small,fill=white,inner sep=1pt]{\(b\)};
\fill (-0.8,0) circle (2.2pt) node[below,font=\small,fill=white,inner sep=1pt]{\(c\)};
\fill (-4,-4) circle (2.2pt) node[below left,font=\small]{\(d\)};
\fill (0,{-4*tanh(0.8/1.5)}) circle (2.2pt) node[left,font=\small,fill=white,inner sep=1pt]{\(e\)};
\fill (0.8,0) circle (2.2pt) node[above right,font=\small,fill=white,inner sep=1pt]{\(f\)};
\draw[->,inkgray] (0.9,1.55)--(0.10,1.92);
\node[dimlab] at (1.22,1.42) {\(B_r\)};
\draw[->,inkgray] (-1.65,-0.55)--(-0.86,-0.06);
\node[dimlab] at (-2.05,-0.62) {\(H_c\)};
\node[dimlab,right] at (2.2,4.55) {饱和};
\draw[dashed,inkgray] (4,0)--(4,4);
\draw[dashed,inkgray] (0,4)--(4,4);
\node[dimlab,below right] at (0,0) {\(O\)};
\fill (0,0) circle (1.8pt);
\end{tikzpicture}
\caption{磁滞回线与初始磁化曲线（重绘，依据原讲义 Figure 3.8-4；虚线为初始磁化曲线，实线为回线；模型曲线只示意形状；原图见下）。}
\end{figure}



\subhdr{3.2 回线面积就是磁滞损耗}

沿回线走一圈，磁场对单位体积磁介质做的功为
\[
W=\oint H\dd B \qquad (\mathrm{J/m^3}).
\]
对闭合回线，这个积分等于回线所围的面积；面积越大，每周期单位体积损耗的能量越多。若样品体积为 \(V\)、频率为 \(f\)，则磁滞损耗功率为
\[
P=f\,V\,W.
\]
这解释了为什么变压器、电机要选软磁材料：不是它们「不会被磁化」，而是它们的回线窄、面积小，反复磁化时发热少（讲义 p6）。

\begin{center}
\includegraphics[width=0.42\linewidth]{fig-3.8-orig/fig3.8-4.png}\\[2pt]
{\small\color{inkgray} 原讲义 Figure 3.8-4（裁剪）：磁滞回线 $a\to b\to c\to d\to e\to f\to a$，标出了 $B_s,B_r,H_c$。}
\end{center}

\subhdr{3.3 软磁与硬磁的回线}

\begin{figure}[htbp]
\centering
\begin{tikzpicture}[x=0.85cm,y=0.85cm]
\begin{scope}
\draw[->] (-3.6,0)--(3.6,0) node[right,font=\small]{\(H\)};
\draw[->] (0,-3.6)--(0,3.6) node[above,font=\small]{\(B\)};
\draw[thick,violet,domain=-3:3,smooth,samples=100] plot (\x,{3*tanh((\x+0.35)/1.2)});
\draw[thick,violet,domain=-3:3,smooth,samples=100] plot (\x,{3*tanh((\x-0.35)/1.2)});
\node[dimlab] at (0,-4.25) {(a) 软磁：回线窄，\(H_c\) 小};
\end{scope}
\begin{scope}[xshift=8.0cm]
\draw[->] (-3.6,0)--(3.6,0) node[right,font=\small]{\(H\)};
\draw[->] (0,-3.6)--(0,3.6) node[above,font=\small]{\(B\)};
\draw[thick,petal,domain=-3:3,smooth,samples=100] plot (\x,{3*tanh((\x+1.8)/0.9)});
\draw[thick,petal,domain=-3:3,smooth,samples=100] plot (\x,{3*tanh((\x-1.8)/0.9)});
\node[dimlab] at (0,-4.25) {(b) 硬磁：回线宽，\(H_c\) 大};
\end{scope}
\end{tikzpicture}
\caption{软磁与硬磁的磁滞回线（重绘，依据原讲义 Figure 3.8-5；宽度为示意；原图见下）。}
\end{figure}

\begin{center}
\includegraphics[width=0.42\linewidth]{fig-3.8-orig/fig3.8-5.png}\\[2pt]
{\small\color{inkgray} 原讲义 Figure 3.8-5（裁剪）：左为软磁（小矫顽力），右为硬磁（大矫顽力）。}
\end{center}

%=====================================================================
\hdr{§4}{示波器测量原理：怎样把 \(B\) 和 \(H\) 变成屏幕上的 \(x\) 和 \(y\)}

\subhdr{4.1 测量电路}

样品是一个带两个绕组的铁芯：原边 \(N_1\) 串一个很小的取样电阻 \(R_1\)，副边 \(N_2\) 串 \(R_2\) 与 \(C\)。原边电流 \(i_1\) 在 \(R_1\) 上产生电压 \(U_{R_1}\)，它正比于 \(H\)；副边感应电动势 \(\varepsilon_2\) 正比于 \(\dd B/\dd t\)，经过 \(R_2\)--\(C\) 积分后，电容电压 \(U_C\) 正比于 \(B\)。示波器放在 X-Y 模式：CH1 接 \(U_{R_1}\) 作 X 轴，CH2 接 \(U_C\) 作 Y 轴，屏幕上就直接画出 \(B\)--\(H\) 回线（图 5；原讲义 Figure 3.8-6 见 §5.1 原图）。

\begin{figure}[H]
\centering
\begin{tikzpicture}[x=0.92cm,y=0.92cm]
% primary loop
\draw[thick] (0.6,2.95) -- (0.6,4.7) -- (2.2,4.7) -- (2.2,4.2);
\draw[thick] (2.2,1.0) -- (2.2,0.5) -- (0.6,0.5) -- (0.6,2.25);
\draw[thick] (0.6,2.6) circle (0.35);
\node at (0.6,2.6) {$\sim$};
\node[lab,left] at (0.18,2.6) {$u_1$};
% R1
\draw[comp] (1.0,0.34) rectangle (1.6,0.66);
\node[lab] at (1.3,0.98) {$R_1$};
% CH1
\draw[thick] (1.0,0.34) -- (1.0,0.05);
\draw[thick] (1.6,0.34) -- (1.6,0.05);
\draw[thick] (1.0,0.05) circle (0.07);
\draw[thick] (1.6,0.05) circle (0.07);
\node[lab] at (1.3,-0.30) {CH1 ($x$)};
% core
\draw[thick] (2.2,0.9) -- (2.2,4.3);
\draw[thick] (2.8,0.9) -- (2.8,4.3);
\draw[thick] (2.2,0.9) -- (2.8,0.9);
\draw[thick] (2.2,4.3) -- (2.8,4.3);
% coils
\foreach \y in {3.8,3.0,2.2,1.4} {
  \draw[thick] (2.2,\y+0.4) arc[start angle=90,end angle=270,radius=0.4];
  \draw[thick] (2.8,\y+0.4) arc[start angle=90,end angle=-90,radius=0.4];
}
\node[lab] at (1.58,2.6) {$N_1$};
\node[lab] at (3.42,2.6) {$N_2$};
\node[font=\footnotesize,fill=white,inner sep=1pt] at (2.5,4.45) {铁芯};
% secondary loop
\draw[thick] (2.8,4.2) -- (2.8,4.7) -- (3.8,4.7);
\draw[comp] (3.8,4.54) rectangle (4.6,4.86);
\node[lab] at (4.2,5.15) {$R_2$};
\draw[thick] (4.6,4.7) -- (6.2,4.7) -- (6.2,4.0);
\draw[thick] (5.9,4.0) -- (6.5,4.0);
\draw[thick] (5.9,3.8) -- (6.5,3.8);
\draw[thick] (6.2,3.8) -- (6.2,0.5) -- (2.8,0.5) -- (2.8,1.0);
\node[lab] at (6.55,3.35) {$C$};
% CH2
\draw[thick] (6.2,4.35) -- (7.0,4.35);
\draw[thick] (7.0,4.35) circle (0.07);
\draw[thick] (6.2,3.45) -- (7.0,3.45);
\draw[thick] (7.0,3.45) circle (0.07);
\node[lab,right] at (7.10,3.9) {CH2 ($y$)};
% current arrows
\draw[arr] (1.0,4.7) -- (1.4,4.7);
\node[lab] at (1.2,5.0) {$i_1$};
\draw[arr] (3.4,4.7) -- (3.0,4.7);
\node[lab] at (3.2,5.0) {$i_2$};
\end{tikzpicture}
\caption{测量电路（重绘，依据原讲义 Figure 3.8-6；原图没有标出元件符号与节点名，本图补上 \(R_1,R_2,C,N_1,N_2\) 与 CH1/CH2 的接法；原图见 §5.1）。}
\end{figure}

\subhdr{4.2 \(H\) 的测量：看原边电流}

闭合铁芯中，安培环路定理给出
\[
H=\frac{N_1 i_1}{L},
\]
其中 \(L\) 是磁路平均长度（样品参数里给的 \(L=0.130\ \mathrm{m}\) 或 \(0.075\ \mathrm{m}\)），不是绕组的长度。取样电阻 \(R_1\) 上的电压是
\[
U_{R_1}=i_1R_1,
\]
所以
\[
H=\frac{N_1}{L}\cdot\frac{U_{R_1}}{R_1}
=\frac{N_1U_{R_1}}{LR_1}.
\]
示波器 CH1 的垂直灵敏度为 \(S_x\)（V/div），屏幕上偏离中心的格数为 \(x\)，则 \(U_{R_1}=S_x x\)，于是
\[
H=\frac{N_1S_x}{LR_1}\,x.
\]
\(R_1\) 要小，是为了不让它明显改变原边电流；但它又不能为零，否则没有电压可测。

\subhdr{4.3 \(B\) 的测量：副边感应 + \(R_2C\) 积分}

副边绕组中的感应电动势由法拉第定律给出：
\[
\varepsilon_2=-N_2\frac{\dd\Phi}{\dd t}=-N_2S\frac{\dd B}{\dd t}.
\]
副边回路是 \(R_2\) 与 \(C\) 串联，回路电流为
\[
i_2=\frac{\varepsilon_2}{R_2+\dfrac{1}{\mathrm{j}\omega C}}.
\]
电容上的电压是
\[
U_C=\frac{i_2}{\mathrm{j}\omega C}
=\frac{\varepsilon_2}{1+\mathrm{j}\omega R_2C}.
\]
当
\[
R_2\gg\frac{1}{\omega C}\quad\Longleftrightarrow\quad \omega R_2C\gg1
\]
时，分母中的 1 可以忽略，\(U_C\approx\varepsilon_2/(\mathrm{j}\omega R_2C)\)，把 \(\varepsilon_2=-N_2S\,\dd B/\dd t\) 代入并对时间积分，得到
\[
U_C\approx\frac{N_2S}{R_2C}\,B.
\]
用时间域的话说：\(R_2\) 很大时，副边电流几乎由 \(R_2\) 决定，\(i_2\approx\varepsilon_2/R_2\)，电容充电电流正比于 \(\varepsilon_2\)，所以
\[
U_C=\frac{1}{C}\int i_2\dd t
\approx\frac{1}{R_2C}\int\varepsilon_2\dd t
=\frac{N_2S}{R_2C}B.
\]
这就是「\(R_2C\) 积分器」的名字来源：它把 \(\dd B/\dd t\) 积回 \(B\)。符号是负的（积分器反相），屏幕上回线可能上下颠倒；读 \(B_r,H_c\) 的大小时取绝对值即可，必要时把 CH2 反相。

\subhdr{4.4 从屏幕坐标到 \(B\)、\(H\)}

CH2 的垂直灵敏度为 \(S_y\)（V/div），屏幕纵坐标格数为 \(y\)，则 \(U_C=S_y y\)。代入上式：
\[
B=\frac{CR_2}{N_2S}\,U_C=\frac{CR_2S_y}{N_2S}\,y.
\]
于是四个待测量可以写成
\[
B_r=\frac{CR_2S_y}{N_2S}\,y_r,\qquad
H_c=\frac{N_1S_x}{LR_1}\,x_c,
\]
\[
\mu_r=\frac{B}{\mu_0H}
=\frac{CR_2LR_1S_y}{\mu_0N_2SN_1S_x}\cdot\frac{y}{x}.
\]
其中 \(y_r\) 是回线与 \(B\) 轴交点（\(x=0\)）的纵坐标，\(x_c\) 是回线与 \(H\) 轴交点（\(y=0\)）的横坐标。预习 p2--p3 的数据表就是这四个式子的排版；把 \(S_x,S_y\) 的单位从 mV/div 换成 V/div 时，别忘了 \(10^{-3}\)。

\begin{center}
\begin{tabular}{@{}p{5.4cm}p{8.8cm}@{}}
\toprule
\thead{量} & \thead{由屏幕读数算} \\
\midrule
\(H\) & \(H=\dfrac{N_1S_x}{LR_1}\,x\) \\
\(B\) & \(B=\dfrac{CR_2S_y}{N_2S}\,y\) \\
\(B_r\) & \(B_r=\dfrac{CR_2S_y}{N_2S}\,y_r\) \\
\(H_c\) & \(H_c=\dfrac{N_1S_x}{LR_1}\,x_c\) \\
\(\mu_r\) & \(\mu_r=\dfrac{CR_2LR_1S_y}{\mu_0N_2SN_1S_x}\cdot\dfrac{y}{x}\) \\
\bottomrule
\end{tabular}
\end{center}

\subhdr{4.5 一个数值例子（只示范代法，不是本次实测）}

设某次测量取 \(R_1=5\ \Omega\)、\(R_2=50\ \mathrm{k}\Omega\)、\(C=1.0\ \mu\mathrm{F}\)、\(S_x=0.5\ \mathrm{V/div}\)、\(S_y=0.2\ \mathrm{V/div}\)，样品 1 的 \(L=0.130\ \mathrm{m}\)、\(S=1.24\times10^{-4}\ \mathrm{m^2}\)、\(N_1=N_2=150\)。读到某点 \(x=1.0\ \mathrm{div}\)、\(y=2.0\ \mathrm{div}\)，则
\[
H=\frac{150\times0.5}{0.130\times5}\times1.0=115.4\ \mathrm{A/m},
\]
\[
B=\frac{1.0\times10^{-6}\times5.0\times10^4\times0.2}{150\times1.24\times10^{-4}}\times2.0
=1.075\ \mathrm{T},
\]
\[
\mu_r=\frac{1.075}{4\pi\times10^{-7}\times115.4}=7.4\times10^3.
\]
先检查积分条件：\(\omega R_2C=2\pi\times50\times5.0\times10^4\times1.0\times10^{-6}=15.7\gg1\)，相位误差约 \(\arctan(1/15.7)=3.6^\circ\)，可以接受。若把 \(C\) 换成 \(0.1\ \mu\mathrm{F}\)，\(\omega R_2C=1.57\)，相位误差约 \(32^\circ\)，回线会明显失真——这就是 §4.6 要处理的问题。

\subhdr{4.6 相位失真与满屏}

\textbf{（1）为什么 \(R_2\) 太小时回线会变形。} 上面的近似只在 \(\omega R_2C\gg1\) 时成立。若 \(R_2<1/(\omega C)\)，\(U_C\) 与 \(B\) 之间有明显的相位差：即使样品是线性的（没有磁滞），屏幕上也会出现一个椭圆；对铁磁样品，回线被「拧」成歪斜的形状，\(B_r\)、\(H_c\) 都读不准（见下方原图 3.8-7）。严格地说，
\[
\frac{U_C}{B}
=-\frac{N_2S}{R_2C}\cdot\frac{1}{1+\dfrac{1}{\mathrm{j}\omega R_2C}},
\]
幅度误差因子是 \(1/\sqrt{1+1/(\omega R_2C)^2}\)，相位误差是 \(\arctan[1/(\omega R_2C)]\)。要求幅度误差小于 1\%，需要 \(\omega R_2C\gtrsim7\)；要求相位误差小于 \(1^\circ\)，需要 \(\omega R_2C\gtrsim57\)。所以「\(\gg1\)」在 50 Hz 下通常取 \(\omega R_2C\ge20\)（相位误差 \(<3^\circ\)）比较稳妥。



\textbf{（2）怎么消。} 优先增大 \(R_2\) 或 \(C\)，把 \(\omega R_2C\) 抬上去；同时保证两路都用 DC 耦合、探头补偿一致（否则两通道本身的相位差也会叠进来）；再微调 \(R_2\)、\(C\)，直到回线光滑闭合、左右对称。

\textbf{（3）怎么让回线占满屏幕。} 回线顶点对应 \((\pm H_{\max},\pm B_{\max})\)，屏幕上对应 \((\pm4,\pm4)\) 格时正好占满 8\(\times\)8 网格。操作顺序是：先把 \(R_2,C\) 调到积分条件满足；再增大激励电流到饱和；然后调 \(R_1\)（改变同样电流下的 \(U_{R_1}\)，即 X 幅度）和 \(R_2,C\)（改变 Y 幅度），配合示波器的 \(S_x,S_y\)，让顶点落在 \((\pm4,\pm4)\)；最后固定并记录 \(R_1,R_2,C,S_x,S_y\) 五个值。注意 \(R_1\) 也会影响原边电流，\(R_2,C\) 也会影响积分条件，所以不要指望一次调好，按「先相位、后幅度、再刻度」的顺序迭代。

\begin{center}
\includegraphics[width=0.42\linewidth]{fig-3.8-orig/fig3.8-7.png}\\[2pt]
{\small\color{inkgray} 原讲义 Figure 3.8-7（裁剪）：$R_2<1/(\omega C)$ 时的畸变回线；横轴 $U_x$、纵轴 $U_y$。}
\end{center}

\begin{figure}[htbp]
\centering
\begin{tikzpicture}[x=0.85cm,y=0.85cm]
\draw[step=1,inkgray!35,very thin] (-4,-4) grid (4,4);
\draw[->] (-4.6,0)--(4.6,0) node[right,font=\small]{$x$ (H)};
\draw[->] (0,-4.6)--(0,4.6) node[above,font=\small]{$y$ (B)};
\foreach \k in {-4,-3,-2,-1,1,2,3,4} {
  \node[font=\tiny,inkgray] at (\k,-0.25) {\k};
  \node[font=\tiny,inkgray] at (-0.25,\k) {\k};
}
\draw[thick,violet,domain=-4:4,smooth,samples=120] plot (\x,{4*tanh((\x+0.8)/1.5)});
\draw[thick,violet,domain=-4:4,smooth,samples=120] plot (\x,{4*tanh((\x-0.8)/1.5)});
\fill (4,4) circle (2pt) node[above right,font=\footnotesize]{$(4,4)$};
\fill (-4,-4) circle (2pt) node[below left,font=\footnotesize]{$(-4,-4)$};
\fill (0,0) circle (1.6pt);
\end{tikzpicture}
\caption{满屏时的示波器画面：回线顶点落在 \((\pm4,\pm4)\) 格，中心在 \((0,0)\)。实际屏幕上每一大格还可以再分小格，读数时估到 0.05 格即可。}
\end{figure}

%=====================================================================
\hdr{§5}{进实验室：一步一步怎么做}

本节步骤编号与讲义 p8--p9 的 Experimental Contents and Procedures 对应。每一步先写“做什么”，再写“为什么”，报告里讨论操作时可以直接引用。

\subhdr{5.1 接线（p8 步骤 1）}

按图 5（重绘，§4.1）或下方原讲义 Figure 3.8-6 接线：原边 \(N_1\) 与 \(R_1\) 串联接交流源，副边 \(N_2\) 与 \(R_2\)、\(C\) 串联；\(R_1\) 两端接示波器 CH1，\(C\) 两端接 CH2。仪器面板上用红色箭头标出了接线方向，样品通过改变接线位置来更换。

\begin{center}
\includegraphics[width=0.52\linewidth]{fig-3.8-orig/fig3.8-6.png}\\[2pt]
{\small\color{inkgray} 原讲义 Figure 3.8-6（裁剪）：测量电路。与图 5 对照看；原图没有标元件数值，接线以仪器面板上的红色箭头为准。}
\end{center}

\warnbox{两个通道的“地”是连在一起的。接线时先确认仪器输出端与示波器探头的地在同一参考点，不要把 CH1、CH2 的地夹子随便夹在 \(R_1\) 或 \(C\) 的两端，否则可能把 \(R_1\) 或 \(C\) 短路，读到的波形看着正常、数值却是错的。接完线先请老师检查，再通电（讲义 p8 的 Note）。}

\subhdr{5.2 初始状态（p8 步骤 2）}

把幅度调节旋钮逆时针拧到最小，确认 \(R_1,R_2,C\) 都不是零、且满足实验要求。这样做的原因是：通电瞬间不能让样品直接进入深饱和，也不能让原边电流过大；从最小激励开始，后面才能单调地往上加。

\subhdr{5.3 开机与 X-Y 模式（p8 步骤 3）}

打开示波器和磁滞回线测量仪，频率调到约 50 Hz。按水平控制区的 MENU 把时基改成 X-Y 模式；按 CH1、CH2 的 POSITION 把光点调到屏幕中心，使中心点对应 \((0,0)\)；两个通道的耦合都设为 DC（在触发控制区的 MENU 里选）。

\textbf{为什么用 DC 耦合}：AC 耦合相当于加了一个高通滤波器，会把低频成分滤掉，回线的形状和相位都会变；DC 耦合才能把 \(U_{R_1}\) 与 \(U_C\) 的真实波形送进 X、Y 放大器。

\subhdr{5.4 调到饱和、再调到满屏（p8 步骤 4）}

顺时针缓慢、\textbf{单调}增大激励电流，观察回线从一个小椭圆逐渐长成典型的饱和回线。判断饱和的标志是：继续增大电流时，回线顶点主要沿水平方向（\(H\) 方向）移动，\(B\) 几乎不再增加，顶部和底部变平、变尖。

然后调 CH1、CH2 的 SCALE 旋钮（按一下可以在粗调/细调之间切换），改变 \(S_x,S_y\)，让回线的正、负顶点分别落在 \((4.00,4.00)\) 与 \((-4.00,-4.00)\) 格。讲义 p8 还给了三个“微调”手段：调 \(R_1\) 改变 \(U_x\)（X 幅度），调 \(R_2,C\) 改变 \(U_y\)（Y 幅度），调 \(S_x,S_y\) 改变刻度。若回线相位歪了，调 \(R_2,C\) 消失真。

\callout{调好以后，\(R_1,R_2,C,S_x,S_y\) 五个值就\textbf{固定不动}，并记在数据表上。后面所有 \(B,H,B_r,H_c,\mu_r\) 的换算都用这五个值；中途改一个，前面读的格数就全部作废。}

\subhdr{5.5 退磁（p9 步骤 5）}

把激励电流\textbf{单调、缓慢}地逆时针减到零，直到屏幕上的回线缩成一个点。这个点的位置应该在 \((0,0)\) 附近；如果偏离，说明还有剩磁或示波器零点没对准，需要重新退磁或重新调零。

\textbf{为什么必须先退磁}：基本磁化曲线要求从 \(B=H=0\) 的 O 点出发。如果样品带着剩磁 \(B_r\) 开始，起点就不在原点，测出来的“初始曲线”整体平移，\(\mu_{\max}\) 和饱和点都会偏。退磁的本质是让磁畴在幅度逐渐减小的交变场中反复翻转，最后停在无规则的多畴状态。

\subhdr{5.6 基本磁化曲线（p9 步骤 6）}

从退磁态开始，\textbf{单调增大}激励电流，使回线正顶点的横坐标依次为
\[
x=0.00,\ 0.40,\ 0.80,\ 1.20,\ 1.60,\ 2.00,\ 2.40,\ 3.00,\ 4.00\ \mathrm{div},
\]
每到一个 \(x\)，记录该顶点的纵坐标 \(y\)。这里的 \(y\) 必须取\textbf{顶点}的 \(y\)，不是回线在某个 \(x\) 处的上下两支；顶点才对应“当前 \(H\) 下已经达到的 \(B\)”。

\callout{读数时估到 \(0.05\ \mathrm{div}\)。\(x=0\) 那一行是退磁点，\(y\) 应为 0；如果不是 0，先检查零点，不要直接当成数据。}

\subhdr{5.7 动态磁滞回线（p9 步骤 7）}

当屏幕上是典型的饱和回线、顶点在 \((\pm4,\pm4)\) 时，记录
\[
x=-4.00,\ -3.00,\ -2.00,\ -1.50,\ -1.00,\ -0.50,\ 0.00,\ 0.50,\ 1.00,\ 1.50,\ 2.00,\ 3.00,\ 4.00
\]
处的 \(y\)。讲义没有说明记哪一支；本讲义建议记\textbf{上支}（从 \(+B_s\) 下降到 \(-B_s\) 的那一支），下支用中心对称 \(B(-H)=-B(H)\) 补出来。这样 \(x=0\) 处的上支 \(y\) 就是 \(+B_r\)，\(y=0\) 处的正 \(x\) 就是 \(+H_c\)。如果报告表格给了两列，就上下支都记。

\subhdr{5.8 样品 2 与比较（p9 第 2 部分）}

换样品 2，重复 5.1--5.7。样品 2 的参数是 \(L=0.075\ \mathrm{m}\)、\(S=1.20\times10^{-4}\ \mathrm{m^2}\)、\(N_1=N_2=150\)，\textbf{不要}沿用样品 1 的 \(L,S\)。比较时看三件事：\(B_r\) 与 \(H_c\) 的大小、回线的宽窄、\(\mu_{\max}\) 的高低；再结合 §1.2 的软磁/硬磁讨论。

\subhdr{5.9 记录清单}

\begin{center}
\begin{tabular}{@{}p{4.6cm}p{3.3cm}p{2.2cm}p{5.0cm}@{}}
\toprule
\thead{项目} & \thead{符号} & \thead{单位} & \thead{何时记} \\
\midrule
取样电阻 & \(R_1\) & \(\Omega\) & 步骤 2 检查、步骤 4 固定 \\
积分电阻 & \(R_2\) & \(\mathrm{k}\Omega\) & 步骤 4 固定 \\
积分电容 & \(C\) & \(\mu\mathrm{F}\) & 步骤 4 固定 \\
垂直灵敏度 & \(S_x,S_y\) & V/div & 步骤 4 固定 \\
频率 & \(f\) & Hz & 步骤 3 设定 \\
基本磁化曲线 & \(x,y\) & div & 步骤 6 逐点 \\
回线读数 & \(x,y\) & div & 步骤 7 逐点 \\
\(B_r,H_c\) 坐标 & \(y_r,x_c\) & div & 步骤 7 读出 \\
\bottomrule
\end{tabular}
\end{center}

\subhdr{5.10 操作要点}

\begin{itemize}
\item 接线要牢，尤其是样品接线柱；接触电阻会直接叠加到 \(R_1\) 的电压上。
\item 调饱和、退磁、测基本曲线时都要\textbf{单调}改变激励；中途来回拧会引入小回线，破坏“从 O 出发”的前提。
\item 固定 \(R_1,R_2,C,S_x,S_y\) 后不要再碰；换量程前先把数据记完。
\item 样品和线圈不要用手长时间捂着，温度变化会让 \(B_r,H_c\) 漂移。
\item 发现回线左右不对称、中心漂移，先查零点与接地，再查样品是否被磁化过。
\end{itemize}

%=====================================================================
\hdr{§6}{数据处理：从格数到 \(B\)、\(H\)、\(\mu_r\)}

\subhdr{6.1 样品参数}

\begin{center}
\begin{tabular}{@{}lccc@{}}
\toprule
\thead{参数} & \thead{样品 1} & \thead{样品 2} & \thead{单位} \\
\midrule
磁路平均长度 \(L\) & 0.130 & 0.075 & m \\
截面积 \(S\) & \(1.24\times10^{-4}\) & \(1.20\times10^{-4}\) & \(\mathrm{m^2}\) \\
原边匝数 \(N_1\) & 150 & 150 & 匝 \\
副边匝数 \(N_2\) & 150 & 150 & 匝 \\
\bottomrule
\end{tabular}
\end{center}

\subhdr{6.2 先算 \(H\) 和 \(B\)，再算 \(B_r,H_c,\mu_r\)}

把 §4.4 的公式抄成一张“代入顺序”：

\begin{enumerate}
\item 由 \(x\) 算 \(H=\dfrac{N_1S_x}{LR_1}x\)；
\item 由 \(y\) 算 \(B=\dfrac{CR_2S_y}{N_2S}y\)；
\item 在回线上读 \(x=0\) 处的 \(y_r\)，算 \(B_r=\dfrac{CR_2S_y}{N_2S}y_r\)；
\item 在回线上读 \(y=0\) 处的 \(x_c\)，算 \(H_c=\dfrac{N_1S_x}{LR_1}x_c\)；
\item 对基本磁化曲线的每一点算 \(\mu_r=\dfrac{B}{\mu_0H}\)，或者直接用
\[
\mu_r=\frac{CR_2LR_1S_y}{\mu_0N_2SN_1S_x}\cdot\frac{y}{x}.
\]
\end{enumerate}

第 5 步里的系数 \(K=\dfrac{CR_2LR_1S_y}{\mu_0N_2SN_1S_x}\) 在同一组 \(R_1,R_2,C,S_x,S_y\) 下是常数，所以 \(\mu_r\propto y/x\)。预习 p3 要求画“\(\mu\sim H\) 曲线 using original data \((y/x\sim x)\)”，就是利用这个正比关系：先画 \(y/x\) 随 \(x\) 的变化，形状与 \(\mu_r\) 随 \(H\) 的变化完全一样，最后再把纵轴乘上 \(K\)。

\subhdr{6.3 基本磁化曲线的表怎么填}

预习 p3 的表格是 \(x/\mathrm{div}\)、\(y/\mathrm{div}\)、\(y/x\) 三行。填法：

\begin{itemize}
\item \(x=0\) 那一列是退磁点，\(y=0\)，\(y/x\) 不填（分母为零）；
\item 其余列先填 \(y\)，再算 \(y/x\)；
\item 用 \(H=\dfrac{N_1S_x}{LR_1}x\)、\(B=\dfrac{CR_2S_y}{N_2S}y\) 把每一列换成 \(H,B\)；
\item 以 \(H\) 为横轴、\(B\) 为纵轴画基本磁化曲线；以 \(H\) 为横轴、\(\mu_r\) 为纵轴画 \(\mu\sim H\) 曲线；
\item 在 \(\mu\sim H\) 曲线上找最大值，记下对应的 \(H\) 与 \(\mu_r\)，就是“最大相对磁导率的坐标与数值”。
\end{itemize}

\subhdr{6.4 磁滞回线的图怎么画}

把步骤 7 的 \((x,y)\) 逐点换成 \((H,B)\)，上支用实测，下支用中心对称补；标出 \((0,B_r)\)、\((H_c,0)\)、\((-H_c,0)\)、\((0,-B_r)\)。如果要估回线面积，可在 \(B\)--\(H\) 图上数格子或梯形积分：
\[
W\approx\sum_i \bar H_i\,\Delta B_i,
\]
单位是 \(\mathrm{J/m^3}\)。本实验不要求算损耗，但报告讨论“软磁为什么回线窄”时可以用面积做定性依据。

\subhdr{6.5 误差从哪里来}

\begin{itemize}
\item \textbf{格数读数}：人眼估读 \(\pm0.05\ \mathrm{div}\)，在回线陡峭处会被放大；多读几次取平均。
\item \textbf{\(S_x,S_y\) 的校准}：示波器垂直增益有 2\%--3\% 的误差；同一通道要用同一档位。
\item \textbf{\(R_1,R_2,C\) 的标称值}：电阻、电容有 5\%--10\% 的容差；若实验室给了实测值，优先用实测值。
\item \textbf{积分器相位}：\(\omega R_2C\) 不够大时，\(U_C\) 与 \(B\) 有相位差，回线“变胖”；用 §4.6 的条件检查。
\item \textbf{频率漂移}：\(\omega\) 变化会改变积分条件和涡流损耗；实验过程中频率应稳定在 50 Hz。
\item \textbf{温度与磁历史}：样品发热、上次测量的剩磁都会让 \(B_r,H_c\) 漂移；测前退磁、测中不要长时间通大电流。
\item \textbf{磁路不均匀}：\(H=N_1i_1/L\) 假设磁场均匀分布在平均磁路上；样品有接缝、气隙或绕组不均匀时，这个假设只是近似。
\end{itemize}

%=====================================================================
\hdr{§7}{报告素材：目的、器材与 10 道预习问题（中英对照）}

以下先给实验目的与器材，再逐题给“英文答案 + 中文对照”。英文可直接抄进报告，中文帮助理解。

\subhdr{7.0 实验目的 / Experimental Objectives}

\begin{enumerate}
\item Understand the fundamental concepts of magnetization curve, hysteresis loop, permeability, coercive force, and residual magnetism.\\
中文：理解磁化曲线、磁滞回线、磁导率、矫顽力和剩磁的基本概念。
\item Learn how to observe the magnetization curve and hysteresis loop of a ferromagnetic material using an oscilloscope.\\
中文：学会用示波器观察铁磁材料的磁化曲线和磁滞回线。
\item Measure permeability, saturation flux density, coercive force, and residual magnetism.\\
中文：测量磁导率、饱和磁感应强度、矫顽力和剩磁。
\item Compare the hysteresis loops under different frequencies for different ferromagnetic materials.\\
中文：比较不同频率、不同铁磁材料的磁滞回线。
\end{enumerate}

\subhdr{7.0b 实验器材 / Experimental Apparatus}

\begin{enumerate}
\item DH4516N apparatus for measuring magnetic hysteresis loops (AC excitation source, R1, R2, C, two samples).\\
中文：DH4516N 磁滞回线测量仪（含交流激励源、R1、R2、C 与两个样品）。
\item Digital oscilloscope with two channels and X-Y mode.\\
中文：数字示波器（双通道，支持 X-Y 模式）。
\item Connecting wires and the printed data table.\\
中文：连接导线与自行打印的数据表。
\item Sample 1: \(L=0.130\ \mathrm{m}\), \(S=1.24\times10^{-4}\ \mathrm{m^2}\), \(N_1=N_2=150\).\\
中文：样品 1：\(L=0.130\ \mathrm{m}\)，\(S=1.24\times10^{-4}\ \mathrm{m^2}\)，\(N_1=N_2=150\)。
\item Sample 2: \(L=0.075\ \mathrm{m}\), \(S=1.20\times10^{-4}\ \mathrm{m^2}\), \(N_1=N_2=150\).\\
中文：样品 2：\(L=0.075\ \mathrm{m}\)，\(S=1.20\times10^{-4}\ \mathrm{m^2}\)，\(N_1=N_2=150\)。
\end{enumerate}

\subhdr{Q1. 为什么铁磁材料的 \(B\) 与 \(H\) 不是线性关系？}
因为磁化来自磁畴：畴壁移动（先可逆、后不可逆）和畴转动对 \(H\) 的响应不是线性的，\(M\) 先快后慢并最终饱和到 \(M_s\)；\(B=\mu_0(H+M)\)，所以 \(B\)--\(H\) 先缓、后陡、再平，\(\mu=B/H\) 不是常数。

\textbf{Answer (EN):} Ferromagnetism comes from magnetic domains. Domain-wall motion and domain rotation respond nonlinearly to H: M rises quickly, then slows and saturates at Ms. Since B = mu0(H + M), mu = B/H is not constant; at high H, B is about mu0(H + Ms) and the curve flattens.

\subhdr{Q2. 说明磁化曲线与磁滞回线，并画图。}
磁化曲线是从退磁态单调增大 \(H\) 得到的初始 \(B\)--\(H\) 曲线，终点是饱和 \(B_s\)。磁滞回线是循环磁化下的闭合路径 O--a--b--c--d--e--f--a：a 是正向饱和，b 是 \(H=0\) 时的剩磁 \(B_r\)，c 是 \(B=0\) 时的矫顽力 \(H_c\)，d 是反向饱和，e、f 是反向的 \(B_r\)、\(H_c\)。回线面积等于每周期单位体积的磁滞损耗。

\textbf{Answer (EN):} The magnetization curve is the initial B-H curve from a demagnetized state: B rises slowly, then steeply, then bends to saturation Bs. The hysteresis loop is the closed path O-a-b-c-d-e-f-a: a is positive saturation, b is Br at H = 0, c is Hc at B = 0, d is negative saturation, e and f are the reverse Br and Hc. Its area equals the hysteresis loss per unit volume per cycle.

\subhdr{Q3. 为什么要测磁化曲线和磁滞回线？}
为了拿到 \(B_s,B_r,H_c,\mu_{\max}\) 这些特征参数：它们决定材料是软磁还是硬磁、适合作变压器铁芯还是永磁体，也用来估算磁滞损耗、检验材料批次是否合格。

\textbf{Answer (EN):} The measurements give Bs, Br, Hc and mur,max. They classify soft/hard materials, help select cores for transformers, motors, relays and permanent magnets, estimate hysteresis loss, and check material quality.

\subhdr{Q4. 示波器测磁滞回线的原理、电路和公式。}
原边 \(N_1\) 的电流 \(i_1\) 在 \(R_1\) 上产生 \(U_{R_1}=i_1R_1\)，所以 \(H=N_1U_{R_1}/(LR_1)\)；副边 \(N_2\) 感应 \(\varepsilon_2=-N_2S\,\dd B/\dd t\)，经 \(R_2C\) 积分（\(\omega R_2C\gg1\)）后 \(U_C=N_2SB/(R_2C)\)。示波器 X-Y 模式：CH1 接 \(U_{R_1}\)（X，\(H\)），CH2 接 \(U_C\)（Y，\(B\)），屏上就是 \(B\)--\(H\) 回线。换算公式见 §4.4。

\textbf{Answer (EN):} The primary current i1 gives UR1 = i1R1, so H = N1UR1/(L R1). The secondary emf epsilon2 = -N2S dB/dt is integrated by R2C (omega R2 C >> 1), giving UC = N2S B/(R2C). In X-Y mode CH1 (UR1) drives X (H) and CH2 (UC) drives Y (B), so the screen shows the B-H loop. With UR1 = Sx x and UC = Sy y: B = C R2 Sy y/(N2S), H = N1 Sx x/(L R1), Br = C R2 Sy yr/(N2S), Hc = N1 Sx xc/(L R1), mur = C R2 L R1 Sy/(mu0 N2 S N1 Sx) * (y/x).

\subhdr{Q5. 测量前为什么要退磁？怎么退？}
不退磁的话，样品带着剩磁 \(B_r\) 开始，基本磁化曲线的起点不在 O 点，测出的曲线整体平移。退磁方法是：在交流激励下把激励电流（幅度）缓慢、单调地减到零，回线逐渐缩小，最后缩成 \((0,0)\) 处的一个点。

\textbf{Answer (EN):} Demagnetization is needed because the B-H curve depends on magnetic history; residual Br shifts the starting point away from O and biases Br, Hc and mumax. With AC excitation on, reduce the excitation current slowly and monotonically to zero until the loop collapses to a point (AC demagnetization).

\subhdr{Q6. 怎样让样品达到饱和？}
单调增大激励电流，观察回线顶点：当继续增大电流时，顶点主要沿 \(H\) 方向移动、\(B\) 几乎不再增加，回线顶部变平、变尖，就是饱和了。此时再增大电流只会让 \(H\) 更大，\(B\) 基本不变。

\textbf{Answer (EN):} Increase the excitation current monotonically and watch the loop tips. Saturation is reached when further increase mainly moves the tip along H while B hardly increases; the loop top/bottom flattens and the tips become sharp. Then M goes to Ms and only mu0H still grows.

\subhdr{Q7. 怎样避免相位失真？}
保证 \(R_2\gg1/(\omega C)\)，即 \(\omega R_2C\gg1\)；否则 \(U_C\) 与 \(B\) 有相位差，回线会变成歪斜的椭圆。做法是增大 \(R_2\) 或 \(C\)，两通道都用 DC 耦合、探头补偿一致，再微调 \(R_2,C\) 直到回线光滑闭合。50 Hz 下建议 \(\omega R_2C\ge20\)。

\textbf{Answer (EN):} The RC integrator needs R2 >> 1/(omega C), i.e. omega R2 C >> 1. Otherwise UC is phase-shifted relative to B and the loop is distorted (a linear core shows an ellipse). Increase R2 or C, use DC coupling on both channels with matched probes, then fine-tune R2/C until the loop is smooth and closed. At 50 Hz, omega R2 C >= 20 is safe.

\subhdr{Q8. 饱和且不失真时，怎样让回线占满屏幕？}
先把 \(R_2,C\) 调到积分条件满足、回线不失真；再增大激励到饱和；然后调 \(R_1\)（X 幅度）、\(R_2,C\)（Y 幅度）和 \(S_x,S_y\)，让回线中心在 \((0,0)\)、顶点落在 \((\pm4,\pm4)\) 格。最后固定并记录 \(R_1,R_2,C,S_x,S_y\)。

\textbf{Answer (EN):} First set R2/C for good integration, then increase the excitation to saturation. Adjust R1 (X amplitude, H), R2/C (Y amplitude, B) and the vertical scales Sx, Sy so the loop is centered at (0,0) and its tips fall at (+4,+4) and (-4,-4) divisions, filling the 8x8 grid without clipping. Keep R1, R2, C, Sx, Sy fixed and record them.

\subhdr{Q9. 怎样测基本磁化曲线？}
先退磁，固定 \(S_x,S_y,R_1,R_2,C\)；单调增大激励，使回线正顶点的 \(x\) 依次为 0、0.40、0.80、1.20、1.60、2.00、2.40、3.00、4.00 格，记录每个顶点的 \(y\)；再用 \(H=N_1S_x x/(LR_1)\)、\(B=CR_2S_y y/(N_2S)\) 换算并作图。

\textbf{Answer (EN):} Demagnetize first and fix Sx, Sy, R1, R2, C. Increase the excitation monotonically so the positive vertex x is 0, 0.40, 0.80, 1.20, 1.60, 2.00, 2.40, 3.00, 4.00 div; record y at each step. Convert with H = N1 Sx x/(L R1) and B = C R2 Sy y/(N2S), then plot B vs H (and y/x vs x for mur).

\subhdr{Q10. 怎样得到相对磁导率？}
逐点算 \(\mu_r=B/(\mu_0H)\)；用屏幕坐标就是 \(\mu_r=CR_2LR_1S_y/(\mu_0N_2SN_1S_x)\cdot(y/x)\)。在同一组参数下 \(\mu_r\propto y/x\)，所以画 \(y/x\) 随 \(x\) 的曲线，形状就是 \(\mu_r\)--\(H\) 曲线；曲线最大值对应 \(\mu_{\max}\)，其横坐标给出对应的 \(H\)。

%=====================================================================

\textbf{Answer (EN):} At each point mur = B/(mu0 H). In screen coordinates mur = C R2 L R1 Sy/(mu0 N2 S N1 Sx) * (y/x). Plot mur vs H (or y/x vs x); the maximum gives mur,max and its coordinate gives the corresponding H.
\hdr{§8}{课后思考题（中英对照）}

\subhdr{1. 为什么测量前必须退磁？}
因为磁化曲线和回线都与磁化历史有关；不退磁，样品带着剩磁，初始点不是 O，\(B_r,H_c,\mu_{\max}\) 都会偏。

\textbf{Answer (EN):} The B-H curve depends on magnetic history; residual Br shifts the start away from O and biases Br, Hc and mumax. Demagnetize by slowly and monotonically reducing the AC excitation current to zero until the loop collapses to a point.

\subhdr{2. 怎样判断样品已经饱和？}
继续增大激励时，回线顶点主要沿 \(H\) 方向移动，\(B\) 几乎不变，顶部变平、变尖；此时 \(M\to M_s\)，再增大 \(H\) 只有 \(\mu_0H\) 的贡献。

\textbf{Answer (EN):} Further increase of the excitation mainly extends H while B hardly changes; the loop tips become sharp and the top/bottom flatten. This means M goes to Ms, so only mu0H still grows.

\subhdr{3. 为什么要求 \(R_2\gg1/(\omega C)\)？不满足会怎样？}
只有满足这个条件，副边电流才近似为 \(\varepsilon_2/R_2\)，电容电压才正比于 \(\int\varepsilon_2\dd t\propto B\)。不满足时 \(U_C\) 与 \(B\) 有相位差，线性样品也会显示成椭圆，铁磁样品的回线被扭曲，\(B_r,H_c\) 读不准。

\textbf{Answer (EN):} This condition makes the secondary current i2 about epsilon2/R2, so UC is about integral(epsilon2 dt) proportional to B. If not met, UC is phase-shifted relative to B; even a linear core shows an ellipse and a ferromagnetic loop is distorted, so Br and Hc are misread.

\subhdr{4. 为什么回线面积代表磁滞损耗？单位是什么？}
沿回线一圈，磁场对单位体积介质做功 \(W=\oint H\dd B\)，对闭合路径等于回线面积；单位是 \(\mathrm{J/m^3}\)。乘样品体积 \(V\) 得每周期能量，再乘频率 \(f\) 得损耗功率。

\textbf{Answer (EN):} Per cycle the field does work W = integral(H dB) per unit volume; for a closed loop this equals the enclosed area. The unit is J/m3; multiplying by sample volume V and frequency f gives the hysteresis power loss.

\subhdr{5. 软磁与硬磁的区别和应用？}
软磁：\(H_c\) 小、\(\mu\) 高、回线窄，易磁化易退磁，用于变压器、继电器、电机铁芯；硬磁：\(H_c\) 大、\(B_r\) 高、回线宽，用于永磁体（钕铁硼、钐钴、铁氧体）。

\textbf{Answer (EN):} Soft magnets: small Hc, high mu, narrow loop, easy to magnetize/demagnetize, used for transformer cores, relays and motor cores. Hard magnets: large Hc and Br, wide loop, used for permanent magnets such as NdFeB, SmCo and ferrite.

\subhdr{6. 为什么基本磁化曲线要从退磁态单调增大 \(H\) 来测？}
这样每个点都落在同一条初始磁化曲线上；如果中途来回改变 \(H\)，会走小回线，测到的 \(B\) 不是“当前 \(H\) 对应的初始值”。

\textbf{Answer (EN):} This keeps every point on the same initial magnetization curve. If H is increased and decreased repeatedly, the sample follows minor loops and the measured B is no longer the initial value for that H.

\subhdr{7. 主要误差来源有哪些？}
格数估读、\(S_x,S_y\) 校准、\(R,C\) 容差、积分器相位、频率漂移、温度与磁历史、磁路不均匀；减小办法是多次读数、用实测元件值、保证 \(\omega R_2C\gg1\)、测前退磁、保持频率稳定。

\textbf{Answer (EN):} Grid reading, Sx/Sy calibration, R and C tolerances, integrator phase error, frequency drift, temperature/magnetic history and non-uniform magnetization. Reduce them by repeated readings, measured component values, omega R2 C >> 1, demagnetizing before each run, and stable frequency/connections.

\subhdr{8. 为什么用 X-Y 模式而不是普通时域波形？}
普通时域是 \(U\)--\(t\)，两个通道各画各的，不能直接给出 \(B\) 与 \(H\) 的函数关系；X-Y 模式把 \(U_{R_1}\)（\(H\)）当横轴、\(U_C\)（\(B\)）当纵轴，逐点合成，屏上直接就是 \(B\)--\(H\) 回线。

\textbf{Answer (EN):} The normal time-base shows U vs t for each channel separately, so it does not directly give B as a function of H. X-Y mode plots UR1 (H) on the horizontal axis and UC (B) on the vertical axis, so the screen directly shows the B-H loop.

\subhdr{8.9 实验体会 / Reflection（约 50 词）}

\textbf{Option A (48 words).} This experiment made the B-H loop concrete: the oscilloscope turned abstract domains into a visible curve. I learned why demagnetization and a good RC integrator matter, and how Br, Hc and mur come from screen coordinates. Careful scaling and stable connections strongly affect the results, so patience matters.\\
中文：这次实验让 B-H 回线变得具体：示波器把抽象的磁畴变成了一条可见的曲线。我明白了退磁和良好的 RC 积分器为什么重要，也学会了如何从屏幕坐标得到 Br、Hc 和 mur。仔细定标和稳定接线会显著影响结果，因此耐心很关键。

\textbf{Option B (51 words).} Measuring the hysteresis loop taught me that a clean result depends on preparation: demagnetization, a correct RC integrator, and fixed scales. I was surprised how strongly R2 and C change the loop shape. The experiment connected domain theory, circuits and data processing, and gave me a concrete feel for magnetic materials.\\
中文：测量磁滞回线让我认识到，干净的结果取决于准备：退磁、正确的 RC 积分器和固定的量程。R2 与 C 对回线形状的影响之大令我意外。这个实验把磁畴理论、电路和数据处理连了起来，让我对磁性材料有了具体的感受。

%=====================================================================
\hdr{§9}{易错点（每条都能指回某一步）}

\begin{itemize}
\item \textbf{不退磁就测基本曲线}：起点带 \(B_r\)，曲线平移；对应 §5.5、§5.6。
\item \textbf{先调 \(S_x,S_y\) 再调 \(R_2,C\)}：相位一变，刚调好的满屏又歪了；顺序应是“先相位、后幅度、再刻度”（§4.6、§5.4）。
\item \textbf{\(R_2\) 太小}：回线变成歪椭圆，\(B_r\) 偏大、\(H_c\) 偏小；对应 §4.6。
\item \textbf{读错支}：在 \(x=0\) 处读了下支，得到 \(-B_r\)；在 \(y=0\) 处读了负 \(x\)，得到 \(-H_c\)。取大小或按上下支分别读（§5.7）。
\item \textbf{单位没换}：\(S_x,S_y\) 给的是 mV/div，\(C\) 给的是 \(\mu\mathrm{F}\)，\(R_2\) 给的是 \(\mathrm{k}\Omega\)；代入前先换成 V/div、F、\(\Omega\)（§4.4）。
\item \textbf{把 \(L\) 当成样品“长度”随便量}：\(L\) 是磁路平均长度，样品 1 与样品 2 不同，不能混用（§6.1）。
\item \textbf{换样品不换参数}：样品 2 的 \(L,S\) 与样品 1 不同，公式里的 \(L,S\) 要跟着换（§5.8）。
\item \textbf{用手捂样品/长时间通大电流}：温度升高，\(B_r,H_c\) 漂移；测完一组让样品冷一下（§6.5）。
\item \textbf{两个通道地线夹错}：可能短路 \(R_1\) 或 \(C\)，波形看着对、数值全错（§5.1 的 warnbox）。
\item \textbf{把 \(\mu=B/H\) 当成切线斜率}：报告用割线值；切线斜率是 \(\dd B/\dd H\)，两者不同（§2.1、G3）。
\end{itemize}

%=====================================================================
\hdr{§10}{缺口清单：两份材料没明说、本讲义补了什么}

\gapitem{推导缺口}{讲义 p7 只写了 \(i_2=\varepsilon_2/\sqrt{R_2^2+(1/\omega C)^2}\) 和 \(U_C=N_2SB/(CR_2)\)，没有写相位、没有写复阻抗推导。本讲义 §4.3 补了 \(U_C=\varepsilon_2/(1+\mathrm{j}\omega R_2C)\) 的复阻抗推导，并给出幅度误差与相位误差公式。}
\gapitem{适用条件}{讲义 p7 说“If \(R_2\gg1/\omega C\)”，但没有给“到底多大才够”。本讲义 §4.6 给出：幅度误差 1\% 需 \(\omega R_2C\gtrsim7\)，相位误差 \(1^\circ\) 需 \(\gtrsim57\)，50 Hz 下建议 \(\ge20\)。}
\gapitem{术语混用}{讲义 p3 说“the slope of the magnetization curve ... gives the permeability”，把割线 \(B/H\) 与切线 \(\dd B/\dd H\) 混在一起。本讲义 §2.1 明确区分，并说明报告按 \(B/(\mu_0H)\) 算。}
\gapitem{操作顺序}{讲义 p8 步骤 4 说可以调 \(R_1\)、\(R_2,C\)、\(S_x,S_y\)，但没有说先调哪个、它们怎样互相影响。本讲义 §4.6 与 §5.4 给出“先相位、后幅度、再刻度”的迭代顺序。}
\gapitem{机制缺口}{讲义 p2 说磁畴“互相抵消”，但没有解释磁畴为什么形成。本讲义 §1.1 用交换相互作用与静磁能的一句话补上（标为教学补充，不是讲义原文）。}
\gapitem{公式来源}{讲义 p2--p3 直接给 \(B_r\)、\(H_c\)、\(\mu_r\) 的公式，没有从 \(U_C\)、\(U_{R_1}\) 推。本讲义 §4.2--§4.4 逐条推了。}
\gapitem{图缺}{讲义 p7 的 Figure 3.8-6 没有标 \(R_1,R_2,C,N_1,N_2\) 与 CH1/CH2 的接法，也没有画两个通道的地。本讲义 §4.1 重绘、§5.1 附原图，补上标签与接法（图注标明“重绘”）。}
\gapitem{读数分支}{讲义 p9 步骤 7 说“record the Y coordinates ... when X coordinates are ...”，没有说记上支还是下支。本讲义 §5.7 明确记上支、下支用中心对称补，并指出 \(B_r\)、\(H_c\) 的读法。}
\gapitem{定标}{讲义 p8 步骤 4 要求顶点到 \((\pm4,\pm4)\)，但没有解释为什么是 4。本讲义 §4.6 说明屏幕是 8\(\times\)8 格，\(\pm4\) 就是满屏且不削顶。}
\gapitem{单位}{讲义 p2 附录给 \(1\ \mathrm{T}=1\ \mathrm{V}\cdot\mathrm{F}\cdot\Omega/\mathrm{m^2}\)，没有推导。本讲义 §0.1 用 \(F\cdot\Omega=s\) 补了单位链。}

%=====================================================================
\hdr{§11}{三句话总结}

\begin{enumerate}
\item 铁磁性的非线性来自磁畴：畴壁移动（可逆/不可逆）+ 畴转动 + 饱和；所以 \(B=\mu_0(H+M)\) 里的 \(M\) 不是 \(H\) 的线性函数。
\item 磁滞回线把 \(B_r\)、\(H_c\)、\(B_s\) 和损耗装进一条闭合曲线；回线面积 \(W=\oint H\dd B\) 就是每周期单位体积的磁滞损耗。
\item 示波器测量的核心是两件事：\(R_1\) 把 \(i_1\) 变成正比于 \(H\) 的电压，\(R_2C\) 积分把 \(\varepsilon_2\) 变成正比于 \(B\) 的电压；只要 \(\omega R_2C\gg1\)、五个参数固定，\(B_r,H_c,\mu_r\) 都能从格数算出来。
\end{enumerate}

%=====================================================================
\hdr{§12}{与两份材料的页码对照}

\begin{center}
\begin{tabular}{@{}p{2.4cm}p{10.6cm}p{3.2cm}@{}}
\toprule
\thead{材料页} & \thead{内容} & \thead{本讲义} \\
\midrule
讲义 p1 & 铁磁材料介绍、实验目的、仪器、原理开头 & §0--§1 \\
讲义 p2 & 磁畴、磁化曲线、\(B=\mu_0\mu_rH\) & §1--§2 \\
讲义 p3 & 空气/铁/钢的曲线、磁导率、\(\mu_r\) 曲线 & §2 \\
讲义 p4 & 剩磁、矫顽力、软磁/硬磁 & §2--§3 \\
讲义 p5 & 磁滞回线 O--a--b--c--d--e--f--a & §3 \\
讲义 p6 & 回线闭合、软硬磁回线、测量原理开头 & §3--§4 \\
讲义 p7 & 测量电路、\(H\)、\(B\)、\(R_2C\) 积分、失真 & §4 \\
讲义 p8 & 接线、初始状态、X-Y、调饱和与满屏 & §5.1--§5.4 \\
讲义 p9 & 退磁、基本磁化曲线、动态回线、样品 2 & §5.5--§5.8 \\
预习 p1 & 实验目标 + 10 道预习问题 & §7 \\
预习 p2 & 样品 2 的 \(B_r,H_c\) 数据表、附录参数 & §6 \\
预习 p3 & 样品 1 的 \(B_r,H_c\)、基本磁化曲线、\(\mu_r\) 表 & §6 \\
\bottomrule
\end{tabular}
\end{center}

\flowerrule

\noindent{\small\color{inkgray} 本讲义依据的两份 PDF 均为英文；正文中的中文解释、推导与“为什么”由本讲义补写。凡标“示意”的曲线图只表示形状，不是本次实测数据；实验后请用实测的 \(R_1,R_2,C,S_x,S_y\) 重新计算。}

\end{document}
